\section{Experimental Evaluation}
\label{sec:eval}
\subsection{Testbed, Scope, and Evaluation Flow}
\label{sec:protocol}\label{sec:static}
The experiments ask whether the platform (Q1) preserves write constraints
across heterogeneous protocols, (Q2) separates prevention from readback and
attribution, and (Q3) closes a team-level objective through traceable,
multi-node interventions. We freeze the integrated archive
$B=\texttt{20260920094610-to4}$ from the September 20, 2026 consolidated
report, recording seed 42, commit \texttt{1585948}, harness hash
\texttt{b203077a581aaba1} over nine benchmark-harness files, and Node v24.19.0
on Windows x64. The command is \texttt{node bench/pipeline.mjs --profile
integrated --seed 42 --cl-seeds 3}. Its immediately preceding same-version run
is reported separately rather than selected for its endpoint, and the ablation
and protocol/fault drill come from separate archives, not pooled with $B$.
Numerical controller and mission
observations use separate REST history queries that request one-second buckets
and average up to six and five latest returned bucket means respectively; the
archive preserves aggregates, not raw buckets.

The testbed runs actual Modbus TCP, OPC UA, MQTT, HTTP, and Modbus
RTU-over-TCP transports against simulated plants, including a cast-film
physics model, a strategy-based film line, and a biaxially oriented
polyethylene terephthalate (BOPET) line. The integrated tool harness is
deterministic \texttt{opencode}; mission leads and workers are scripted
mock-engine members, and LLM decisions are disabled. Thus, the experiments
evaluate supervisory execution and its evidence trail, not learned
optimization, physical production performance, or an advantage over another
agent framework. The simulator-derived reference and author-defined probes
are internal oracles, not independent industrial benchmarks.

Because the loop is the object under test, Fig.~\ref{fig:chanloop} states what
an iteration consists of before any number is reported: which signals return
from the plant, where admission is decided, which node executes a write over
which transport, and what evidence survives the write. The measured results
then attach to those stages in a fixed order --- coverage and protocol
behavior (Sec.~\ref{sec:suite}), prevention versus readback
(Sec.~\ref{sec:e1lite}), fixed-controller loop closure
(Sec.~\ref{sec:clbench}), and the recorded team-level mission
(Sec.~\ref{sec:agentloop}) --- so that a later claim can be traced back to the
stage that produced it. Each subsection names its archive and states what the
run does \emph{not} establish, so the section reads as a sequence of
independent checks rather than one aggregate score.

\begin{figure*}[!t]
\centering
\includegraphics[width=0.88\textwidth]{figures/publication/fig3-channel-loop.png}
\caption{The agent-team loop that the evaluation exercises, from the
observation that starts an iteration to the compensation that can end one.
The channel coordinates a lead and workers whose per-node bindings gate
admission---bound scope, availability, approval policy, hard range plus the
active recipe window---before a driver write with register readback where the
protocol supports it. \texttt{dcw\_judge} records the agent's verdict but
neither computes it nor actuates; compensation is a separate, rejectable
write; the human authority and the auto backstop sit outside the channel.
Numbered items are loop steps, not the subsections of Sec.~\ref{sec:mech};
the plant is simulated.}
\label{fig:chanloop}
\end{figure*}

\subsection{Integrated Coverage and Protocol Results}
\label{sec:suite}\label{sec:port}\label{sec:dynamic}\label{sec:plc}\label{sec:govsurface}
Archive $B$ reports 75 passes, zero warnings, and zero failures across 17
phases. A pass is a check verdict, not an independent trial or a statistical
confidence statement. The checks partition by function as 9 plant and
provisioning, 5 protocol I/O and write governance, 7 three-cycle tool loops,
12 records, approval, recipes, and parameters, 16 team objectives and the
BOPET line, 22 controller, scenario reuse, and backstop, and 4 multimodal
acquisition and services. Acquisition-only nodes have no write endpoint; three
omitted write/governance sub-checks must not be counted as tested writes. The
optional LLM phase is absent. A separate source audit
(\texttt{20260920095700-1e3s}) passes three checks, including 20/20
code-constant anchors, and the accompanying \texttt{\_r7-apilive.log} records
60 API passes and two persistence-related skips despite the consolidated
report's zero-skip summary; these overlapping layers are not summed as
independent experiments.

Five acquisition nodes store 59 samples. Four writable stacks each reject six
out-of-constraint probes, giving 24/24 rejections and zero false blocks on
four legal writes; the checker counts HTTP error responses to configured
hard-range and window probes, so it does not verify each physical non-effect
or establish adversarial robustness. Reprovisioning the second film-line
scenario requires no driver-code changes: its three writable lines reject 9/9
probes without false blocks, alongside two acquisition-only satellites.

\begin{table}[!t]
\centering
\caption{Protocol measurements in $B$. Latency quantiles use six timed writes per stack: p50 is the fourth sorted value and p95 the maximum. Readback is one sampled deviation, not an accuracy bound. A dash means unmeasured or inapplicable.}
\label{tab:plc}
\footnotesize
\setlength{\tabcolsep}{2.9pt}
\begin{tabular}{@{}lrrr@{}}
\toprule
Stack & Samples & p50/p95 (ms) & Readback \\
\midrule
Modbus TCP & 9 & 49.232/101.164 & $0.04\,^\circ$C \\
OPC UA & 11 & 19.437/22.110 & 0 rpm \\
MQTT & 12 & 19.427/23.084 & --- \\
HTTP & 13 & 32.941/33.546 & --- \\
RTU over TCP & 14 & --- & --- \\
\bottomrule
\end{tabular}
\end{table}

Table~\ref{tab:plc} retains the per-protocol measurements instead of reducing
integration quality to one score. The report's 30.259-ms aggregate is the
arithmetic mean of the four protocol p50 values, not the median of pooled
writes. MQTT and HTTP archive no measured register readback, so their
transport acceptance must not be described as verified actuator state. The
approval request-to-decision latency is 448\,ms with a scripted approver;
human response time is not measured. The backstop rollback verdict is first
observed at 130.205\,s, and baseline restoration is verified subsequently,
after an additional 1.5-s wait and read; the reported interval is therefore
verdict-observation latency, not measured restoration time or a hard recovery
deadline.

\subsection{Ablating Prevention and Readback}
\label{sec:e1lite}
The separate four-arm archive \texttt{20260920094442-e1lite} uses seed 42 and
three repetitions per arm. These are benchmark-only switches requiring
\texttt{AW\_BENCH\_MODE=1}, not normal operator controls: the ablation
bypasses soft-limit enforcement, readback-error enforcement, or both while
retaining the engineering-range guard, and only the recipe soft layer is
populated in this fixture. The report label ``No readback'' therefore means
the tolerance is relaxed, not that readback I/O stops, and because the mock
driver returns the submitted value without checking tolerance the switch is
inert for mismatch detection here. The experiment isolates window rejection,
not the effectiveness of readback verification. Table~\ref{tab:e1lite} shows
all arms, including the executed violations: each repetition runs six
out-of-constraint probes and evaluates three legal and three boundary cases
separately, on new fixture objects from the same design. Arm-specific seeds
vary numerical probes within the same case families, so these are not paired
identical treatments or independent physical-plant trials. The arms keep the
table order: 1~full, 2~no interlock, 3~no readback, and 4~both disabled.

\begin{table}[!t]
\centering
\caption{Four-arm governance ablation (three repetitions, 18 out-of-constraint probes per arm). ``Boundary'' counts decisions matching the arm-specific boundary oracle; ``breaches'' counts executed window violations. The engineering guard remains active in every arm.}
\label{tab:e1lite}
\footnotesize
\setlength{\tabcolsep}{3pt}
\begin{tabular}{@{}lrrrr@{}}
\toprule
Arm & Rejected & Breaches & False blocks & Boundary \\
\midrule
Full & 18/18 & 0 & 0 & 9/9 \\
No interlock & 12/18 & 6 & 0 & 9/9 \\
No readback & 18/18 & 0 & 0 & 9/9 \\
Both disabled & 12/18 & 6 & 0 & 9/9 \\
\bottomrule
\end{tabular}
\end{table}

Removing the interlock admits both window-class violations in each repetition,
whereas bypassing readback verification does not change rejection
(Table~\ref{tab:e1lite}). This supports a specific mechanistic claim: the
window check prevents these violations; readback and journaling do not
substitute for it. The fourth arm's ``Ungated'' label is not guard-free:
engineering limits remain enforced. All nine legal writes
per arm are accepted. Boundary agreement means the expected accept/reject
decision, not universal acceptance: full and no-readback reject six of nine
boundary cases, whereas the two interlock-disabled arms accept all nine.
Per-repetition p50 ranges are 110.6--141.0, 114.8--127.2, 111.9--127.1, and
124.8--139.9\,ms in table order; each quantile uses 23 successful writes
(three legal and 20 timed), and arms run in fixed order. These descriptive
ranges use mock in-memory readback, including in the verification-bypassed
arms; they do not establish latency equivalence or the cost of removing
physical readback I/O.

A second drill tests the same distinction on the live link. The archive
\texttt{20260920094335-1b5g}
passes 13 checks (three source checks, six API checks, and four protocol
checks); all five driver tests succeed. Its Modbus drill accepts a write in
25\,ms, reads back a setpoint of 182, and observes the simulated process reach
179.5 after 30\,s; five additional violations are rejected. After a
frozen-process injection and subsequent write, the checker first observes a
matching alarm at a rounded 2 s from its post-write timer, not 2 s from fault
injection. A TCP disconnect/reconnect drill detects and recovers the
connection. These narrowly exercised failures support the distinction between
an accepted command and a responding process, not general fault tolerance.

\subsection{Fixed-Controller Loop Closure}
\label{sec:clbench}
To isolate the execution path from agent strategy, the cast-film controller
starts at screw speed $N=150$\,rpm, line speed $v=95$\,m/min, and zone
setpoint $200\,^\circ$C. It maximizes
\begin{equation}
J=55J_{\rm thick}+25J_{\rm quality}+8J_{\rm energy}+7J_{\rm throughput},
\label{eq:objective}
\end{equation}
subject to melt temperature in $[195,225]\,^\circ$C and pressure at most
22\,MPa. For thickness $h$ in $\mu$m, defect percentage $d$ in percentage
points, and commanded speeds $N$ and $v$ in rpm and m/min, respectively,
\begin{equation}
\begin{aligned}
J_{\rm thick}&=\max\!\left\{0,1-\frac{\max\{|h-50|-2,0\}}{10}\right\},\\
J_{\rm quality}&=1-\frac{\min\{d,8\}}{8},\\
J_{\rm energy}&=1-\frac{N-50}{150},\qquad
J_{\rm throughput}=\frac{v}{120}.
\end{aligned}
\label{eq:reward-components}
\end{equation}
On the configured nonnegative-defect and command domain, $J\in[0,95]$;
infeasible states are unscored. The denominator is the stored, three-decimal
reference score, not a rescore of its separately rounded output variables.
The fixed rule updates $N\leftarrow N(50/h)$ with a 22-rpm step limit, holds
line speed fixed, and corrects heating outside $[196,224]\,^\circ$C. The
separately maintained simulator provides a finite-grid reference
$J^*=89.894$, not a proven global optimum: its grid ties the three zone
temperatures and varies temperature from 195 to 225\,$^\circ$C in
2.5\,$^\circ$C steps, screw speed from 50 to 200\,rpm in 5-rpm steps, and line
speed from 20 to 120\,m/min in 2.5-m/min steps, at a fixed 1.0-mm die gap.
Unlike the online controller, the reference can vary line speed; the ratio
therefore compares different action sets, not two equally equipped optimizers.
Energy and throughput use issued setpoints rather than independent sensor
measurements.

\begin{table}[!t]
\centering
\caption{Fixed-controller results in $B$. $J_{\rm end}$ is the mean of the two post-write evaluations, not the last observation. All seeds use two accepted writes and zero rejections.}
\label{tab:controller}
\footnotesize
\setlength{\tabcolsep}{4pt}
\begin{tabular}{@{}rrrrr@{}}
\toprule
Seed & $J_0$ & $J_{\rm end}$ & $J_{\rm end}/J^*$ & Final $h$ ($\mu$m) \\
\midrule
42 & 70.167 & 86.845 & 0.966 & 49.200 \\
43 & 68.608 & 87.236 & 0.970 & 49.683 \\
44 & 65.740 & 87.080 & 0.969 & 50.100 \\
\bottomrule
\end{tabular}
\end{table}

The reported mean ratio is 0.968 and mean wall time is 36.178\,s per seed.
Table~\ref{tab:controller} distinguishes the report's endpoint statistic from
the actual observations: the second write lowers $J$ in every seed, so the
trajectory is not monotone. Final thickness meets the reward band, and the
residual gap to the reference lies in the quality, energy, and throughput
terms rather than in thickness. Because the controller inverts the simulator's
own $h\propto N/v$ relation, this experiment shows loop closure under a known
model, not general search efficiency or learned control superiority.

\subsection{The Recorded AgentTeam Optimization Process}
\label{sec:agentloop}
The first mission links task delegation, a five-minute acquisition query, a
governed temperature write from 200 to $204.704\,^\circ$C, a
$204.700\,^\circ$C observation, an attributed \texttt{keep} record, and
parent-task completion. All six checks pass; the $\pm0.75\,^\circ$C tolerance
makes this a workflow-closure test rather than a precision-control result.

The BOPET mission extends that workflow to nine simulated devices with 30
writable setpoints and 19 acquisition nodes, against a thickness objective of
$25.0\pm0.7\,\mu$m. Figure~\ref{fig:agentteam} reconstructs the actual
intervention sequence from $B$ rather than an illustrative optimization curve.
Four of six allowed writes act on three distinct parameters: casting speed
from 32 to 33.1\,m/min, fast-roll speed from 118 to 120.8\,m/min, outlet rail
width from 3000 to 3050\,mm, then casting speed again to approximately
33.7\,m/min. The archived thickness observations are 26.90, 26.82, 26.46,
26.00, and 25.48\,$\mu$m, so the endpoint sits 0.48\,$\mu$m from the target
and 0.22\,$\mu$m inside the permitted band.

\begin{figure*}[!tb]
\centering
\includegraphics[width=0.94\textwidth]{figures/tii-final/fig-agentteam-mission.pdf}
\caption{Recorded scripted AgentTeam execution in $B$. \textbf{(a)}~The five
archived thickness observations against the $25.0\pm0.7\,\mu$m objective;
markers are returned-bucket means from queries that requested one-second
aggregation, and raw buckets and timestamps are not retained, so no continuous
trajectory is inferred. \textbf{(b)}~The four governed writes, each named by
parameter, setpoint change, and transport. Every write opened an optimization
record that received a logged \texttt{keep}, and the parent task reached
\texttt{COMPLETED}. The plant is simulated and the policy is scripted, not an
LLM-generated search.}
\label{fig:agentteam}
\end{figure*}

Each accepted write creates a distinct optimization record, followed by a
settling gauge observation and a scripted \texttt{keep} judgment. The policy
uses relative-thickness-error gains of 0.5, 0.35, and 0.3 with fixed actuator
rotation, so it does not discover the sequence. The mission logs no
independent post-write register value for these actuators and the caller
issues \texttt{keep} unconditionally, so the demonstrated contribution is
scoped, attributable multi-node execution toward one objective rather than an
independently validated judgment or an optimal recipe window.
Melt-temperature observations remain between 291.1 and 291.525\,$^\circ$C,
inside the configured 268--300\,$^\circ$C window; no disturbance tests that
bound.

\subsection{Repeatability and Limits of the Evidence}
\label{sec:repeatability}
The immediately preceding same-seed, same-hash archive
$A=\texttt{20260920093330-md4}$ also passes 75/75 checks, rejects 24/24
primary and 9/9 second-scenario probes, accepts six controller writes, and
closes the BOPET mission with four writes on three parameters. Their final
thicknesses nevertheless differ, at 25.64 versus 25.48\,$\mu$m, as do mean
controller ratios (0.970 versus 0.968) and backstop verdict-observation
latencies (130.182 versus 130.205\,s). These are two consecutive regression
runs, not a reliability campaign. The comparator reports matching verdicts,
but its controller counter rows contain null fields and its protocol arrays
are empty, so we check the underlying run records rather than treat its
``REPRODUCIBLE'' label as proof of full numerical agreement.

Earlier archives include process collapse despite a non-failing hard gate;
those different-revision runs are not pooled with $A$ and $B$ or erased by
their success. Initial simulator state, sampling timing, and shared-instance
residue remain confounders, and the biax journal check has an empty evidence
array although the mission log contains a casting-node journal excerpt. These
limitations motivate pinned initial conditions, explicit missing-evidence
failures, broader fault injection, and physical validation before deployment.
